


% ich bin ein kommentar

% latex ignoriert whitespace, leerzeichen, absätze etc.
% ein neuer absatz entsteht durch eine leerzeile (mehrere absätze)
% mehrere leerzeilen: trotzdem nur ein absatz
% ein einfacher zeilenumbruch ohne neuen absatz macht man mit \\

\chapter{Grundlagen}

Mit \% schreiben wir einen Kommentar. \\
Latex ignoriert Whitespace, Leerzeichen, Absätze etc.\\
Mehrere Leerzeilen: trotzdem nur ein Absatz.\\
Einen einfachen Zeilenumbruch ohne neuen Absatz macht man 
mit \textbackslash\textbackslash .\\

\section{Gliederung}
Wir gliedern Text mit chapter, section, subsection, subsubsection und paragraph.
Normalerweise sollte man nur bis subsection benutzen, sonst das Dokument irgendwie 
weiter aufteilen. Die letzten beiden eher nicht benutzen. \\ Wenn man im Inhaltsverzeichnis
einen anderen Namen will als in der Kapitelüberschrift, wird der Name fürs
Inhaltsverzeichnis in eckigen Klammern vor die geschweiften Klammern gesetzt, also so: \\
\textbackslash section [Name fürs Inhaltsverzeichnis] \{Name für die Kapitelüberschrift\}

\section{Zeichen}

\subsection{Escapen}
Mit \textbackslash       escapen wir Zeichen wie \% . \\
Einen Backslash machen wir mit \textbackslash textbackslash. \\ 
Zeichen, die escaped werden müssen, sind: \% , \$ , \{\} , \textbackslash,
\& , \# , \_ , \^ und \~ . \\

\subsection{Anführungszeichen}

%Text in normalen "Anführungszeichen" sieht so aus. \\
%Typographisch richtig werden die ``Anführungszeichen'' so.
%Mit jeweils zwei \`` Backticks und \' Apostrophen. \\
Deutscher Text in Anführungszeichen: "`Anführungszeichen"' mit 
jeweils Anführungszeichen und Backtick vorne bzw. Anführungszeichen 
und Apostroph hinten. \\


\subsection{Bindestriche}


Bindestrich kurz - \\
Gedankenstrich lang mit zwei Bindestrichen -- \\
noch länger mit drei Bindestrichen --- \\

\section{Formatierungen}

\subsection{Zentrierung und Absätze}


\begin{center}
    \large
    Dieser Text ist zentriert. \\
    \tiny
    mit begin und end center
\end{center} 


\noindent mit noindent: Absatz nicht eingerückt. \\

\textbackslash newpage macht eine neue Seite

%\newpage

\subsection{Schriftgrößen}


Verschiedene Schriftgrößen sind relativ, nicht absolut.
Es gibt normalsize und kleinere und größere Schriften. \\
Kleiner: tiny, scriptsize, footnotesize, small. \\
Größer: large, Large, LARGE, huge, HUGE. \\


\noindent Beispiele folgen: \\
\begin{tiny} This text is tiny. \end{tiny} \\
\begin{scriptsize} This text is scriptsize. \end{scriptsize} \\
\begin{footnotesize} This text is footnotesize. \end{footnotesize} \\
\begin{small} This text is small. \end{small} \\
\begin{normalsize} This text is normalsize. \end{normalsize} \\
\begin{large} This text is large. \end{large} \\
\begin{Large} This text is Large. \end{Large} \\
\begin{LARGE} This text is LARGE. \end{LARGE} \\
\begin{huge} This text is huge. \end{huge} \\
\begin{Huge} This text is Huge. \end{Huge} \\



\subsection{Schriftauszeichnung}


\noindent 
\textit{textit steht für italic, 
der Text ist dann das Argument in geschweiften Klammern}\\
\textsl{textsl steht für slanted, also etwas schräg}\\
\textbf{textbf steht für boldface}\\
\texttt{texttt steht für typewriter, also wie Schreibmaschine}\\
\textsc{textsc steht für smallcaps, also Kapitälchen}\\
\emph{emph heißt: emphasised, also kursiv und betont im Gegensatz zu textit nur kursiv}\\

Direkte Formatierung ist der Notnagel, nicht der Normalfall. Semantische 
Bezeichnungen sind besser.


\section{Semantik}

Mit Kommandos, LaTeX-Befehlen, kann ich bestimmte Formatierungen einzelnen Klassen
zuordnen und somit semantisch nutzen und leicht alle auf einmal verändern. \\
Beispiel: \\
\Name{Eigennamen} möchte ich immer fett gedruckt haben, \Begriff{Fachbegriffe} immer kursiv. \\
Auch Abkürzungen für häufig gebrauchte Wörter wie zB \LG sind machbar. \\

\section{Komplexe Inhaltselemente}
\subsection{Punktliste}
\subsection{Geordnete Liste}
\subsection{Definitionsliste}
\subsection{Geschachtelte Liste}
\subsection{Tabellen}
\subsection{Abbildungen}
\subsection{Mathematische Formeln}
\subsection{URLs}